\Article{3}
If in place of the discrete points $M,M',M''$ etc. we take a continuous body filling a region $S$ with uniform density $=C$, the sum
\[
M\Phi(m,M)+M'\Phi(m,M')+M''\Phi(m,M'')+\ \text{etc.}
\]
becomes the integral $C\int\mathrm{d}S.\Phi(m,\mathrm{d}S)$ taken over the entire region $S$. By analogy we denote by $(m,\mathrm{d}S)$ the distance from point $m$ to any element of the region $S$, denoted by $\mathrm{d}S$.

If in addition, in place of the discrete points $m,m',m''$ etc., we consider a continuous body filling a region $s$ with uniform density $=c$, the calculation of $\Omega$ will require a double integration. It is to be carried out by first finding, for an unspecified point $\mu$, the value of the expression
\SourceLeaf{42}
\[
-gz+\tfrac12c\int\mathrm{d}s.\varphi(\mu,\mathrm{d}s)+C\int\mathrm{d}S.\Phi(\mu,\mathrm{d}S)
\]
where $z$ is the height of point $\mu$ above plane $H$, and the first integral is to be taken over the entire region $s$, the second over the entire region $S$. This value depends only on the position of point $\mu$. If we denote it by $[\mu]$, then
\[
\Omega=c\int\mathrm{d}s.[\mathrm{d}s]
\]
with the integration extended over the entire region $s$.

This may be written more briefly as follows:
\[
\Omega=-gc\int z\mathrm{d}s+\tfrac12cc\int\int\mathrm{d}s.\mathrm{d}s'.\varphi(\mathrm{d}s,\mathrm{d}s')+cC\int\int\mathrm{d}s.\mathrm{d}S.\Phi(\mathrm{d}s,\mathrm{d}S)
\]
Here $s,s'$ properly denote one and the same region, filled by the movable body, but it must be resolved into its elements twice for the double integration.

\Article{4}
What distinguishes fluid bodies is the perfect mobility of even their smallest parts, so that they can assume any shape and yield to even the slightest force tending to change it. In fluids that do not expand, that is, liquids, to which our investigation is devoted, the volume of every particle must remain constant through all changes of shape. Let us therefore consider a fluid body whose motion is limited by a fixed solid body, a vessel. In addition to gravity, we assume that the particles experience both their mutual attraction and the attraction of the parts of the vessel. Equilibrium will require that the value of $\Omega$ be a maximum, that is, that no infinitely small rearrangement of the parts of the fluid can give $\Omega$ a positive increment. Now the value of $\Omega$ can clearly change only insofar as the shape of the region filled by the whole fluid changes. Motion within the fluid alone cannot change it. There will therefore be equilibrium whenever $\Omega$ cannot increase under any infinitely small change of that shape compatible with the shape of the vessel, with the volume remaining constant. It follows directly that if the shape cannot change at all, the vessel surrounding and touching the fluid on every side, those forces acting on the fluid cannot produce internal motion in it and instead balance one another.

\Article{5}
We proceed to a more detailed investigation of the expression $\Omega$, which must be regarded as the foundation of the theory of fluid equilibrium. Starting \SourceLeaf{43}with the first term, it is immediately clear that $\int z\mathrm{d}s$ represents the product of the volume of the region $s$ and the height of its center of gravity above the plane $H$. The products of that same height with the mass and with the weight of the fluid are accordingly $c\int z\mathrm{d}s$ and $gc\int z\mathrm{d}s$, respectively. If the parts of the fluid were subject to no force other than gravity, the height of the center of gravity would have to be as small as possible in equilibrium. From this it is readily inferred that the free part of the surface, or its free parts, must lie in one and the same horizontal plane, bounding the fluid from above.

\Article{6}
The expansion of the second and third terms concerns two particular cases of a general problem. Given any two regions, each element of the first is to be paired with each element of the second, and the products of three factors are to be added together. The factors are the volume of the element of the first region, the volume of the element of the second, and a given function of the distance between them. The second term concerns the case in which the two regions are identical, and the third the case in which one region is entirely outside the other. The full problem comprises two more cases, where one region is a part of the other or where they share a common part. Although the first two cases could suffice for our purpose and the other two could readily be reduced to them, it will be worthwhile to take up the problem in full generality, since it is of considerable interest in its own right. In this general investigation we shall denote the regions by $s,S$ and the function of distance by the symbol $\varphi$. In applying the result to the second term we must substitute for $S$ the region $s$ itself, and in applying it to the third we must substitute for the function $\varphi$ the function $\Phi$ itself. The integral in question is therefore
\[
\int\int\mathrm{d}s.\mathrm{d}S.\varphi(\mathrm{d}s,\mathrm{d}S)
\]
which has the appearance of a double integration. In fact, since the elements of each of the two regions depend on three variables, it involves a sixfold integration, and we shall now show how to reduce it to a fourfold integration.

